% concatenated from Smatrix4.tex
\documentclass[a4paper,11pt]{article}

\usepackage{amsmath,amssymb,amsfonts,amsthm}
\usepackage[dvipdfmx]{graphicx}
\numberwithin{equation}{section}

%%%%%%%%%%%%%%%%%%%%%%%%%%%%%%%%%%%%%%%%%%%%%%%%%%%%%%%%%%%%
% Basic macros
%%%%%%%%%%%%%%%%%%%%%%%%%%%%%%%%%%%%%%%%%%%%%%%%%%%%%%%%%%%%
\newcommand{\R}{\mathbb{R}}
\newcommand{\C}{\mathbb{C}}
\newcommand{\Ss}{\mathbb{S}}

\newcommand{\dom}{\mathcal{D}}
\newcommand{\Ran}{\operatorname{Ran}}
\newcommand{\diag}{\operatorname{diag}}
\newcommand{\slim}{\operatorname*{s-lim}}

\newcommand{\jj}[1]{j_{#1}}
\newcommand{\hj}[1]{h^{(1)}_{#1}}

\renewcommand{\Im}{\operatorname{Im}}
\renewcommand{\Re}{\operatorname{Re}}

%%%%%%%%%%%%%%%%%%%%%%%%%%%%%%%%%%%%%%%%%%%%%%%%%%%%%%%%%%%%
% theorem environments
%%%%%%%%%%%%%%%%%%%%%%%%%%%%%%%%%%%%%%%%%%%%%%%%%%%%%%%%%%%%
\newtheorem{thm}{Theorem}[section]
\newtheorem{lem}[thm]{Lemma}
\newtheorem{cor}[thm]{Corollary}
\newtheorem{prop}[thm]{Proposition}

\theoremstyle{definition}
\newtheorem{defn}[thm]{Definition}

\theoremstyle{remark}
\newtheorem{rem}[thm]{Remark}

%%%%%%%%%%%%%%%%%%%%%%%%%%%%%%%%%%%%%%%%%%%%%%%%%%%%%%%%%%%%

\title{Determinant Formulas for Scattering Matrices of Schr\"odinger Operators with Finitely Many Concentric $\delta$-Shells}
\author{Masahiro Kaminaga}
\date{}

\begin{document}
\maketitle

\noindent\textbf{Keywords:} $\delta$-interaction; scattering matrix; resolvent formula; partial-wave decomposition

\noindent\textbf{Mathematical Subject Classification (2020):} 35J10, 35P25, 47A40, 81U20

\begin{abstract}
We study stationary scattering for Schr\"odinger operators in $\R^3$ with
finitely many concentric $\delta$--shell interactions of constant real
strengths. Starting from the self--adjoint realization and the boundary
resolvent formula for this model, we show that, after partial--wave
reduction, the same finite-dimensional boundary matrices that arise in the
resolvent formula also determine the channel scattering coefficients.
More precisely, for each angular momentum $\ell$, the channel coefficient
$S_\ell(k)$ satisfies
$S_\ell(k)=\det K_\ell(k^2-i0)/\det K_\ell(k^2+i0)$ for almost every $k>0$,
where $K_\ell(z)=I_N+m_\ell(z)\Theta$ is the $\ell$--th reduced boundary
matrix. Thus, in each channel, the positive--energy scattering problem is
reduced to a finite-dimensional matrix problem, and the scattering phase is
recovered from $\det K_\ell(k^2+i0)$.

We then study the first nontrivial case of two concentric shells in the
$s$--wave channel, where the interaction between the shells produces
nontrivial threshold effects. We derive an explicit formula for $S_0(k)$
and analyze its behavior as $k\downarrow0$. In the regular threshold
regime, we obtain an explicit scattering length. We further identify a
threshold--critical configuration characterized by the existence of a
nontrivial zero--energy radial solution, regular at the origin, whose
exterior constant term vanishes. In the corresponding nondegenerate
exceptional case, the usual finite scattering length breaks down, and
instead $S_0(k)\to -1$ as $k\downarrow0$.
\end{abstract}

\section{Introduction}

Schr\"odinger operators with singular interactions are a classical source of
explicit models in spectral and scattering theory, 
as discussed in \cite{AlbeverioEtAl2005} and in
the work of Brasche, Exner, Kuperin, and \v{S}eba \cite{BrascheEtAl1994}. 
Among them, concentric
$\delta$--shell interactions provide a natural radial model. Since the
interaction is supported on concentric spheres, the operator is rotationally
symmetric and the scattering problem admits a partial--wave decomposition. At
the same time, when two or more shells are present, the interaction among the
shells produces nontrivial scattering and threshold effects. This makes the
model simple enough for explicit calculation, but still rich enough to show
interesting phenomena.

In this paper we study stationary scattering for
\begin{equation}\label{eq:intro-H}
H
=
-\Delta+\sum_{j=1}^N \alpha_j\delta(|x|-R_j)
\qquad \text{in } L^2(\R^3),
\end{equation}
where $0<R_1<\cdots<R_N$ and
$\alpha_1,\ldots,\alpha_N\in\R$, and we compare $H$ with the free operator
$H_0=-\Delta$.

For finitely many concentric spherical $\delta$--shells, Shabani
\cite{Shabani1988} studied the model by self--adjoint extension methods and
reduced it to radial one--dimensional equations with matching conditions in
each partial wave. On the scattering side, Hounkonnou, Hounkpe, and Shabani
\cite{HounkonnouHounkpeShabani1997} studied scattering theory for finitely
many sphere interactions supported by concentric spheres. Related spectral
and scattering properties for radially symmetric penetrable wall models were
studied by Ikebe and Shimada \cite{IkebeShimada1991}. These penetrable wall
models may be viewed as a regular counterpart of spherical shell
interactions. 
Thus, in the earlier literature, the scattering problem is
treated after partial--wave reduction by solving radial equations and
imposing matching conditions at the shells in each channel. 
In particular, the existence of a partial--wave description for this model is not new.

Our starting point is the boundary resolvent framework obtained in \cite{KaminagaAAMP2026}. 
For finitely many concentric shells, that paper
gives a self--adjoint realization of $H$ and a resolvent formula in terms of
a boundary operator. This framework is closely related to the general theory
of hypersurface $\delta$--interactions; see, for example,
\cite{BehrndtLangerLotoreichik2013}. It is also related to abstract
Kre\u{\i}n--type resolvent formulas for singular perturbations; see
\cite{Posilicano2001}. 
In addition, the construction in
\cite{KaminagaAAMP2026} allows the surface strengths $\alpha_j$ to be bounded
real--valued functions on the shells and does not assume that they are constant. 
In the present paper we do not repeat that construction. Instead,
we restrict ourselves to the rotationally symmetric case of constant shell
strengths and study the positive--energy scattering problem from the boundary
resolvent formula. For the convenience of the reader, in
Section~\ref{sec:model} we recall only the precise input from
\cite{KaminagaAAMP2026} that is needed below, namely the quadratic--form
realization of $H$, the interface description of $D(H)$, the boundary
resolvent formula, the trace--class property of the resolvent difference, and
the characterization of the point spectrum by the boundary matrix $K_N(z)$.

The main point of the present paper is that, in the rotationally symmetric
case, the same reduced boundary matrices that arise in the resolvent formula
also determine the channel scattering coefficients. For each angular
momentum $\ell\ge0$, let $S_\ell(k)$ denote the channel scattering
coefficient. Our main result is the determinant formula
\begin{equation}\label{eq:intro-main}
S_\ell(k)
=
\frac{\det K_\ell(k^2-i0)}{\det K_\ell(k^2+i0)},
\qquad \text{for almost every } k>0.
\end{equation}
Here $K_\ell(z)$ is the $\ell$--th reduced boundary matrix. The full boundary
operator has the form
$$
K_N(z)=I+m(z)\Theta,
\qquad
\Theta=\diag(\alpha_1R_1^2,\ldots,\alpha_NR_N^2),
$$
where $m(z)$ is the boundary operator associated with the free Green
function and $I$ is the identity operator on the corresponding boundary
space. Its $\ell$--th partial--wave component is
$$
K_{N,\ell}(z)=I_N+m_\ell(z)\Theta,
$$
where $I_N$ denotes the $N\times N$ identity matrix. To simplify notation, we
write $K_\ell(z)$ for $K_{N,\ell}(z)$. 
Thus \eqref{eq:intro-main} is not merely an explicit channel formula. 
It shows that the positive--energy scattering data are encoded by exactly the same
reduced boundary matrices that appear in the resolvent formula. 
To the best of our knowledge, the determinant representation
\eqref{eq:intro-main}, written explicitly in terms of the reduced boundary
matrix arising from the resolvent formula, does not appear in the earlier
partial--wave literature on concentric $\delta$--shells.

As an application of this framework, we study the first nontrivial case of
two concentric shells. In the $s$--wave channel we derive an explicit formula
for the scattering coefficient and analyze its low--energy behavior as
$k\downarrow0$. In the regular threshold regime, the phase shift admits the
standard expansion
\begin{equation}\label{eq:intro-scatt-length}
\delta_0(k)=-a_{\rm s}k+o(k)
\qquad (k\downarrow0),
\end{equation}
which defines the scattering length $a_{\rm s}$. We obtain an explicit
formula for $a_{\rm s}$ and show that it agrees with the coefficient in the
asymptotics of the zero--energy radial solution, normalized so that, for
$|x|>R_2$,
$$
u(x)=1-\frac{a}{|x|},
$$
where $a=a_{\rm s}$. We also identify a threshold--critical configuration
characterized by the existence of a nontrivial zero--energy radial solution
which is regular at the origin and whose exterior constant term vanishes. In
the corresponding nondegenerate exceptional case, the usual finite
scattering length does not exist and one has
$$
S_0(k)\to -1
\qquad (k\downarrow0).
$$
This gives a concrete zero--energy interpretation of the threshold anomaly
in the double--shell model.

The paper is organized as follows. In Section~\ref{sec:model} we recall the
quadratic form realization and the resolvent formula from
\cite{KaminagaAAMP2026}. In Section~\ref{sec:scatt} we prove the determinant
formula for the channel scattering matrices. In
Section~\ref{sec:double} we specialize to the double $\delta$--shell case and
derive an explicit formula for the $s$--wave scattering coefficient. In
Section~\ref{sec:low} we analyze the low--energy behavior, including the
regular threshold regime and the nondegenerate exceptional threshold regime,
and give a zero--energy interpretation of the latter.
%%%
\section{The model and the resolvent formula}\label{sec:model}

In this section we recall, in the present concentric-shell setting,
the precise ingredients from \cite{KaminagaAAMP2026} that will be used later.
We do not repeat the proofs.

\subsection{The operator and its quadratic form}

For $j=1,\ldots,N$, let
$$
S_j=\{x\in\R^3:|x|=R_j\}.
$$
We consider the Schr\"odinger operator
$$
H=-\Delta+\sum_{j=1}^N \alpha_j\delta(|x|-R_j)
$$
in $L^2(\R^3)$, where $\alpha_1,\ldots,\alpha_N\in\R$.
Schr\"odinger operators with $\delta$--interactions supported on
hypersurfaces are well studied; see, for example,
\cite{AlbeverioEtAl2005,BrascheEtAl1994,BehrndtLangerLotoreichik2013}.
For abstract Kre\u{\i}n--type resolvent formulas for singular perturbations,
see also \cite{Posilicano2001}. For finitely many concentric spherical
shells, a self--adjoint realization and a boundary integral resolvent formula
were obtained in \cite{KaminagaAAMP2026}.

We define the sesquilinear form $h$ on the form domain
$D[h]=H^1(\R^3)$ by
\begin{equation}\label{eq:form}
h[u,v]
=
\int_{\R^3}\nabla u\cdot\nabla\overline{v}\,dx
+
\sum_{j=1}^N \alpha_j\int_{S_j}u\overline{v}\,d\sigma_j,
\qquad
u,v\in H^1(\R^3).
\end{equation}
Here $D[h]$ denotes the form domain of $h$, and $d\sigma_j$ denotes the
surface measure on $S_j$. Since each $S_j$ is a smooth compact hypersurface,
the trace map
\[
H^1(\R^3)\to L^2(S_j)
\]
is bounded. 
Hence, by the trace inequality, the boundary terms are form bounded with
relative bound zero with respect to the Dirichlet form, and thus $h$ is
closed and lower semibounded.

\subsection{Boundary operators and the resolvent formula}

Let
$$
H_0=-\Delta
$$
denote the free Hamiltonian. Let $z\in\C\setminus[0,\infty)$ and choose
$\sqrt{z}$ so that $\Im\sqrt{z}>0$. We write
$$
R_0(z)=(H_0-z)^{-1},
\qquad
G_z(x,y)=\frac{e^{i\sqrt{z}|x-y|}}{4\pi|x-y|}.
$$
Here $R_0(z)$ is the free resolvent and $G_z(x,y)$ is the corresponding free
Green function.

For each $j=1,\ldots,N$, let $\tau_j:H^1(\R^3)\to L^2(S^2)$ denote the trace
on $S_j$ transported to the unit sphere by the parametrization $x=R_j\omega$,
that is,
$$
(\tau_j u)(\omega)=u(R_j\omega),
\qquad
\omega\in S^2.
$$
Thus $\tau_j u$ is the trace of $u$ on $S_j$, viewed as a function on $S^2$.
Accordingly, all surface integrals defining the boundary operators below are
written with respect to $d\omega$ on $S^2$ rather than $d\sigma_j$ on $S_j$.
We define the single--layer operators
$$
(\Gamma_j(z)\varphi)(x)
=
\int_{S^2}G_z(x,R_j\omega)\varphi(\omega)\,d\omega,
\qquad
j=1,\ldots,N,
$$
and define
$$
\Gamma(z):\bigoplus_{j=1}^N L^2(S^2)\to L^2(\R^3)
$$
by
$$
\Gamma(z)(\varphi_1,\ldots,\varphi_N)
=
\sum_{j=1}^N \Gamma_j(z)\varphi_j.
$$
We also introduce the operator matrix
$$
m(z)=\bigl(m_{ij}(z)\bigr)_{i,j=1}^N
$$
on $\bigoplus_{j=1}^N L^2(S^2)$, where each entry
$m_{ij}(z)$ is the operator on $L^2(S^2)$ given by
$$
(m_{ij}(z)\varphi)(\omega)
=
\int_{S^2}G_z(R_i\omega,R_j\omega')\varphi(\omega')\,d\omega'.
$$
Following \cite{KaminagaAAMP2026}, we set
\begin{equation}\label{eq:Theta}
\Theta=\diag(\alpha_1R_1^2,\ldots,\alpha_NR_N^2),
\qquad
K_N(z)=I+m(z)\Theta,
\end{equation}
where $I$ denotes the identity operator on
$L^2(S^2)\oplus\cdots\oplus L^2(S^2)$. Thus $K_N(z)$ is the boundary
operator matrix associated with the shell interaction.

The next theorem collects the only facts from \cite{KaminagaAAMP2026}
that will be used in the present paper.

\begin{thm}\label{thm:input}
The quadratic form $h$ in \eqref{eq:form} is closed and lower semibounded
on $H^1(\R^3)$, and therefore defines a self--adjoint operator $H$ in
$L^2(\R^3)$.

Moreover, a function $u\in L^2(\R^3)$ belongs to the operator domain $D(H)$
if and only if $u$ is piecewise $H^2$ on the regions separated by the shells,
is continuous across each sphere $r=R_j$, and satisfies
\begin{equation}\label{eq:jump-intro-model}
\partial_r u(R_j+0,\omega)-\partial_r u(R_j-0,\omega)
=
\alpha_j u(R_j,\omega),
\qquad
j=1,\ldots,N,
\quad
\omega\in S^2.
\end{equation}
Here $\partial_r u(R_j\pm0,\omega)$ denotes the radial derivative taken from
the exterior and interior sides of the sphere $r=R_j$.

Let $z\in\C\setminus[0,\infty)$. If $K_N(z)$ is invertible, then
\begin{equation}\label{eq:resolvent}
(H-z)^{-1}
=
(H_0-z)^{-1}
-\Gamma(z)\Theta K_N(z)^{-1}\Gamma(\overline{z})^*.
\end{equation}
Here $\Gamma(\overline{z})^*$ denotes the adjoint of $\Gamma(\overline{z})$.
Moreover,
$$
(H-z)^{-1}-(H_0-z)^{-1}
$$
is trace class.

Finally, for $z\in\C\setminus[0,\infty)$, the operator $K_N(z)$ is
noninvertible if and only if $z\in\sigma_{\mathrm p}(H)$, where
$\sigma_{\mathrm p}(H)$ denotes the point spectrum of $H$.
\end{thm}

Theorem~\ref{thm:input} is a specialization of the self--adjoint realization,
the boundary resolvent formula, and the spectral characterization of the
boundary operator proved in \cite{KaminagaAAMP2026}.

In particular, since $H$ is self--adjoint, $i\notin\sigma_{\mathrm p}(H)$,
and hence $K_N(i)$ is invertible. The trace--class resolvent difference in
Theorem~\ref{thm:input} will be used in Section~\ref{sec:scatt} to obtain
the existence and completeness of the wave operators.

\subsection{Partial--wave reduction}

Since $\alpha_1,\ldots,\alpha_N$ are real constants, both $H$ and the boundary
operator $m(z)$ are rotationally symmetric. Accordingly, on each copy of
$L^2(S^2)$, the operator $m(z)$ is diagonal with respect to the spherical
harmonic decomposition, and the same is true for the direct sum
$$
L^2(S^2)\oplus\cdots\oplus L^2(S^2).
$$

Let $\{Y_{\ell m}\}_{\ell\ge0,\,-\ell\le m\le\ell}$ be an orthonormal basis
of spherical harmonics on $S^2$. For each $\ell\ge0$, the restriction of
$m(z)$ to the $\ell$--th spherical harmonic sector is described by an
$N\times N$ complex matrix, which we denote by $m_\ell(z)$. We denote by
$K_{N,\ell}(z)$ the $\ell$--th partial--wave component of the boundary
operator matrix $K_N(z)$ introduced in \eqref{eq:Theta}. For simplicity, we
write $K_\ell(z)$ for $K_{N,\ell}(z)$ in the rest of the paper.

\begin{lem}\label{lem:mell}
Let $z\in\C\setminus[0,\infty)$ and set $k=\sqrt{z}$ with $\Im k>0$.
Then $m_\ell(z)$ is the $N\times N$ matrix whose entries are given by
\begin{equation}\label{eq:mell-entry}
(m_\ell(z))_{ij}
=
ik\,\jj{\ell}\!\bigl(k\min\{R_i,R_j\}\bigr)\,
\hj{\ell}\!\bigl(k\max\{R_i,R_j\}\bigr),
\qquad 1\le i,j\le N,
\end{equation}
where $\jj{\ell}=j_\ell$ is the spherical Bessel function and
$\hj{\ell}=h_\ell^{(1)}$ is the outgoing spherical Hankel function.
We also denote by $h_\ell^{(2)}$ the incoming spherical Hankel function.
In particular, $(m_\ell(z))_{ij}=(m_\ell(z))_{ji}$.
Moreover, the $\ell$--th partial--wave component of $K_N(z)$ is
\begin{equation}\label{eq:Kell}
K_\ell(z)=I_N+m_\ell(z)\Theta,
\end{equation}
where $I_N$ denotes the $N\times N$ identity matrix.
\end{lem}

\begin{proof}
Let
$$
r_<:=\min\{|x|,|y|\},
\qquad
r_>:=\max\{|x|,|y|\},
$$
and write
$$
\widehat x=\frac{x}{|x|},
\qquad
\widehat y=\frac{y}{|y|}.
$$
The spherical harmonic expansion of the free Green function reads
$$
G_z(x,y)
=
ik\sum_{n=0}^\infty\sum_{q=-n}^n
j_n(kr_<)\,h_n^{(1)}(kr_>)\,
Y_{nq}(\widehat x)\overline{Y_{nq}(\widehat y)}.
$$
Hence, for $\varphi=Y_{\ell m}$, we obtain
\begin{eqnarray*}
(m_{ij}(z)Y_{\ell m})(\omega)
&=&
\int_{S^2}G_z(R_i\omega,R_j\omega')Y_{\ell m}(\omega')\,d\omega'
\\
&=&
ik\,j_\ell\!\bigl(k\min\{R_i,R_j\}\bigr)\,
h_\ell^{(1)}\!\bigl(k\max\{R_i,R_j\}\bigr)\,
Y_{\ell m}(\omega),
\end{eqnarray*}
by orthonormality of the spherical harmonics. This proves
\eqref{eq:mell-entry}.

Since $\Theta=\diag(\alpha_1R_1^2,\ldots,\alpha_NR_N^2)$ acts only on the
shell index and does not mix spherical harmonics, the restriction of
$$
K_N(z)=I+m(z)\Theta
$$
to the $\ell$--th partial wave is exactly
$$
K_\ell(z)=I_N+m_\ell(z)\Theta.
$$
This completes the proof.
\end{proof}

Thus, in each angular momentum channel, the boundary operator matrix $K_N(z)$
reduces to the finite matrix $K_\ell(z)$.
%%%%%%%%%%%%%%%%%%%%%%%%%%%%%%%%%%%%%%%%%%%%%%%%
\section{Scattering matrices for finitely many concentric shells}
\label{sec:scatt}

\subsection{Free spectral representation and partial--wave decomposition}

To describe scattering, we work in the standard spectral representation of the
free Hamiltonian $H_0=-\Delta$; see, for example,
\cite[Ch.~XI]{ReedSimonIII} and \cite{Yafaev1992}.

At a fixed energy $k^2>0$, let $u_{\rm in}$ be a fixed free solution of
$$
(-\Delta-k^2)u_{\rm in}=0
\qquad \text{in }\R^3.
$$
We call a solution $u$ of
$$
(H-k^2)u=0
$$
an outgoing solution with incident part $u_{\rm in}$ if $u_{\rm in}$ is a
free solution of
$$
(-\Delta-k^2)u_{\rm in}=0
$$
and $u-u_{\rm in}$ satisfies the Sommerfeld radiation condition
$$
\lim_{r\to\infty} r(\partial_r-ik)(u-u_{\rm in})(x)=0.
$$
In this case, $u_{\rm in}$ is called the incident wave.
Here $r=|x|$ and $\partial_r$ denotes the radial derivative. 
In the present partial--wave setting, we will take the incident part to be
$$
j_\ell(k|x|)Y_{\ell m}(\widehat x),
$$
where $\widehat x=x/|x|$. Since the spherical harmonics form a complete
system on $S^2$ and the problem is rotationally invariant, it suffices to
consider incident waves of this form. The corresponding outgoing term will
be described by outgoing spherical Hankel functions.

Let $\widehat f$ denote the Fourier transform of $f\in L^2(\R^3)$,
$$
\widehat f(\xi)
=
(2\pi)^{-3/2}\int_{\R^3}e^{-ix\cdot\xi}f(x)\,dx.
$$
We define
$$
(\mathcal F_0 f)(k,\omega)=\widehat f(k\omega),
\qquad
k>0,\ \omega\in S^2,
$$
initially for sufficiently regular $f$. Then, by Plancherel's theorem and the spherical change of variables,
$\mathcal F_0$ extends to a unitary operator from $L^2(\R^3)$ onto
$$
L^2\bigl((0,\infty),k^2dk;L^2(S^2,d\omega)\bigr),
$$
where $d\omega$ denotes the standard surface measure on the unit sphere $S^2$.
In this representation, $H_0=-\Delta$ is diagonalized as multiplication by
$k^2$, that is,
$$
(\mathcal F_0 H_0 f)(k,\omega)
=
k^2(\mathcal F_0 f)(k,\omega),
\qquad f\in D(H_0).
$$

For each $k>0$, the fiber space is $L^2(S^2)$, and the spherical harmonic
decomposition gives
$$
L^2(S^2)
=
\bigoplus_{\ell=0}^\infty {\cal H}_\ell,
\qquad
{\cal H}_\ell
=
\operatorname{span}\{Y_{\ell m}:-\ell\le m\le\ell\}.
$$
Equivalently, for each fixed $k>0$,
$$
(\mathcal F_0 f)(k,\omega)
=
\sum_{\ell=0}^\infty\sum_{m=-\ell}^{\ell}
(\mathcal F_0 f)(k,\ell,m)\,Y_{\ell m}(\omega),
$$
where
$$
(\mathcal F_0 f)(k,\ell,m)
=
\int_{S^2}(\mathcal F_0 f)(k,\omega)\,
\overline{Y_{\ell m}(\omega)}\,d\omega.
$$

Since the shell strengths $\alpha_1,\ldots,\alpha_N$ are constants, both $H$
and $H_0$ commute with the natural unitary action of the rotation group on
$L^2(\R^3)$. As a consequence, the scattering operator admits a partial--wave
decomposition in the free spectral representation. 

The next theorem records this partial--wave decomposition.

\begin{thm}\label{thm:wave}
The wave operators
$$
W_\pm(H,H_0)
=
\slim_{t\to\pm\infty} e^{itH}e^{-itH_0}
$$
exist and are complete. The scattering operator
$$
S=W_+(H,H_0)^*W_-(H,H_0)
$$
is unitary on $L^2(\R^3)$.
In the free spectral representation $\mathcal F_0$,
there exists a measurable family of unitary operators $S(k^2)$ on $L^2(S^2)$
such that
$$
(\mathcal F_0 S\mathcal F_0^{-1}g)(k,\omega)=S(k^2)g(k,\omega)
$$
for almost every $k>0$. Moreover, for each $\ell\ge0$ there exists a scalar
function $S_\ell(k)$, defined for almost every $k>0$, such that
\begin{equation}\label{eq:Sdiag}
(\mathcal F_0 S\mathcal F_0^{-1}g)(k,\ell,m)
=
S_\ell(k)\,g(k,\ell,m).
\end{equation}
Equivalently,
$$
S(k^2)\upharpoonright_{{\cal H}_\ell}=S_\ell(k)I_{{\cal H}_\ell}
$$
for almost every $k>0$.
\end{thm}

\begin{proof}
By Theorem~\ref{thm:input},
$$
(H-i)^{-1}-(H_0-i)^{-1}
$$
is trace class. Hence the wave operators exist and are complete by the
Birman--Kuroda theorem; see \cite[Ch.~XI]{ReedSimonIII}.
Since $H_0=-\Delta$ has purely absolutely continuous spectrum, the
scattering operator
$$
S=W_+(H,H_0)^*W_-(H,H_0)
$$
is unitary on $L^2(\R^3)$.

Since both $H$ and $H_0$ commute with the action of the rotation group, the
wave operators and hence $S$ commute with rotations. Moreover, by the
intertwining property of the wave operators and the functional calculus,
$S$ commutes with the spectral measure of $H_0$. Therefore, in the free
spectral representation $\mathcal F_0$, the operator
$\mathcal F_0 S\mathcal F_0^{-1}$ is decomposable with respect to the energy
parameter $k^2$. Thus there exists a measurable family of unitary operators
$S(k^2)$ on $L^2(S^2)$ such that
$$
(\mathcal F_0 S\mathcal F_0^{-1}g)(k,\omega)=S(k^2)g(k,\omega)
$$
for almost every $k>0$.

Since $S$ commutes with rotations, each fiber operator $S(k^2)$ commutes with
the rotation action on $L^2(S^2)$ for almost every $k>0$. Because each
spherical harmonic subspace ${\cal H}_\ell$ is irreducible under rotations,
Schur's lemma implies that
$$
S(k^2)\upharpoonright_{{\cal H}_\ell}=S_\ell(k)I_{{\cal H}_\ell}
$$
for almost every $k>0$. This is equivalent to \eqref{eq:Sdiag}.
\end{proof}

%%%
\subsection{Outgoing solutions in a fixed partial wave}

We now fix a partial wave and construct the corresponding outgoing solutions.

\begin{lem}
\label{lem:Kell-boundary}
For each $\ell\ge0$ and each $k>0$, the boundary values
\[
m_\ell(k^2\pm i0),
\qquad
K_\ell(k^2\pm i0)
\]
exist. Moreover, if
\[
m_\ell^\pm(k):=m_\ell(k^2\pm i0),
\qquad
K_\ell^\pm(k):=K_\ell(k^2\pm i0),
\]
then
\[
m_\ell^-(k)=m_\ell^+(k)^*,
\]
and
\[
\det K_\ell^-(k)=\overline{\det K_\ell^+(k)}.
\]
\end{lem}

\begin{proof}
By Lemma~\ref{lem:mell}, the entries of $m_\ell(z)$ are given explicitly in
terms of spherical Bessel and Hankel functions. These expressions admit
boundary values on $(0,\infty)$, so that $m_\ell(k^2\pm i0)$ exist for every
$k>0$, and hence so do $K_\ell(k^2\pm i0)$.
%
For the upper boundary value, Lemma~\ref{lem:mell} gives
\[
(m_\ell^+(k))_{ij}
=
ik\,\jj{\ell}\bigl(k\min\{R_i,R_j\}\bigr)\,
\hj{\ell}\bigl(k\max\{R_i,R_j\}\bigr).
\]
For the lower boundary value, we approach the cut $(0,\infty)$ from the lower
half-plane while keeping the branch determined by $\Im\sqrt{z}>0$ on
$\C\setminus[0,\infty)$. Hence
\[
\sqrt{k^2-i0}=-k,
\qquad
\sqrt{k^2+i0}=k,
\]
for $k>0$.
Using
\[
\jj{\ell}(-t)=(-1)^\ell\jj{\ell}(t),
\qquad
\hj{\ell}(-t)=(-1)^\ell h_\ell^{(2)}(t),
\qquad t>0,
\]
we obtain
\[
(m_\ell^-(k))_{ij}
=
-ik\,\jj{\ell}\bigl(k\min\{R_i,R_j\}\bigr)\,
h_\ell^{(2)}\bigl(k\max\{R_i,R_j\}\bigr).
\]
Since $\jj{\ell}(t)$ is real for $t>0$ and
\[
\overline{\hj{\ell}(t)}=h_\ell^{(2)}(t),
\qquad t>0,
\]
it follows that
\[
m_\ell^-(k)=m_\ell^+(k)^*.
\]

By definition,
\[
K_\ell^\pm(k)=I_N+m_\ell^\pm(k)\Theta.
\]
Since $\Theta=\Theta^*$ is a real diagonal matrix, Sylvester's determinant
identity
\[
\det(I+AB)=\det(I+BA)
\]
yields
\begin{eqnarray*}
\det K_\ell^-(k)
&=&
\det\bigl(I_N+m_\ell^-(k)\Theta\bigr)
\\
&=&
\det\bigl(I_N+\Theta m_\ell^-(k)\bigr)
\\
&=&
\det\bigl(I_N+\Theta m_\ell^+(k)^*\bigr)
\\
&=&
\det\bigl((I_N+m_\ell^+(k)\Theta)^*\bigr) = \overline{\det K_\ell^+(k)}.
\end{eqnarray*}
This proves the claim.
\end{proof}

We construct the outgoing solution in the $(\ell,m)$ channel corresponding to
a given incident wave.

\begin{lem}
\label{lem:partial-wave-construction}
Fix $\ell\ge0$, fix $m$ with $-\ell\le m\le\ell$, and let
$k>0$ satisfy
\[
\det K_\ell(k^2+i0)\neq0.
\]
Write
\[
K_\ell^+(k):=K_\ell(k^2+i0),
\qquad
b_\ell(k)
=
{}^t\bigl(\jj{\ell}(kR_1),\ldots,\jj{\ell}(kR_N)\bigr),
\]
and define
\[
c_\ell(k):=K_\ell^+(k)^{-1}b_\ell(k)\in\C^N.
\]
Then there exists an outgoing solution in the $(\ell,m)$ channel whose
incident part is
\[
\jj{\ell}(k|x|)Y_{\ell m}(\widehat x).
\]
For $|x|=r>R_N$, this solution has the form
\begin{equation}\label{eq:ext-form-lemma}
u_{\ell m}^+(x,k)
=
\frac12
\Bigl(
h_\ell^{(2)}(kr)+\sigma_\ell(k)\hj{\ell}(kr)
\Bigr)
Y_{\ell m}(\widehat x),
\end{equation}
where
\begin{equation}\label{eq:sigma-lemma}
\sigma_\ell(k)
=
1-2ik\,{}^t b_\ell(k)\Theta K_\ell^+(k)^{-1}b_\ell(k).
\end{equation}
\end{lem}

\begin{proof}
By assumption, the matrix
\[
K_\ell^+(k)=K_\ell(k^2+i0)
\]
is invertible, so $c_\ell(k)$ is well defined.

Let
\[
m_\ell^+(k):=m_\ell(k^2+i0),
\]
and let
\[
G_k^+(x,y):=\frac{e^{ik|x-y|}}{4\pi|x-y|}.
\]
For $j=1,\ldots,N$, define
\[
\Phi_j(x,k)
:=
\int_{S^2} G_k^+(x,R_j\omega')\,Y_{\ell m}(\omega')\,d\omega'.
\]
Equivalently,
\[
\Phi_j(x,k)
=
R_j^{-2}
\int_{S_j} G_k^+(x,y)\,Y_{\ell m}(y/R_j)\,d\sigma_j(y),
\]
where $d\sigma_j$ denotes the surface measure on $S_j$.

Thus $\Phi_j(\cdot,k)$ is smooth on $\R^3\setminus S_j$, satisfies
\[
(-\Delta-k^2)\Phi_j(\cdot,k)=0
\qquad \text{in }\R^3\setminus S_j,
\]
and is continuous across each sphere.

For $i\neq j$, the function $\Phi_j(\cdot,k)$ is smooth across $S_i$, hence
\[
\partial_r\Phi_j(R_i+0,\omega,k)-\partial_r\Phi_j(R_i-0,\omega,k)=0.
\]

For $i=j$, using the explicit formulas for $\Phi_j$ inside and outside $S_j$,
which will be derived below in \eqref{eq:Phi-exterior-lemma} and
\eqref{eq:Phi-interior-lemma}, we obtain, at $r=R_j$,
\begin{align*}
&\partial_r\Phi_j(R_j+0,\omega,k)-\partial_r\Phi_j(R_j-0,\omega,k)
\\
&\qquad
=
ik^2\Bigl(
j_\ell(kR_j)\bigl(h_\ell^{(1)}\bigr)'(kR_j)
-
h_\ell^{(1)}(kR_j)j_\ell'(kR_j)
\Bigr)Y_{\ell m}(\omega).
\end{align*}
Using the Wronskian identity
\[
j_\ell(t)\bigl(h_\ell^{(1)}\bigr)'(t)-j_\ell'(t)h_\ell^{(1)}(t)
=
\frac{i}{t^2},
\]
we obtain
\[
\partial_r\Phi_j(R_j+0,\omega,k)-\partial_r\Phi_j(R_j-0,\omega,k)
=
-R_j^{-2}Y_{\ell m}(\omega).
\]
Therefore
\begin{equation}\label{eq:jump-Phi-lemma}
\partial_r\Phi_j(R_i+0,\omega,k)-\partial_r\Phi_j(R_i-0,\omega,k)
=
-\delta_{ij}R_i^{-2}Y_{\ell m}(\omega).
\end{equation}

We next record two formulas for $\Phi_j$.

First, evaluating $\Phi_j$ on $S_i$ and using Lemma~\ref{lem:mell}, we obtain
\begin{equation}\label{eq:trace-Phi-lemma}
\Phi_j(R_i\omega,k)
=
(m_\ell^+(k))_{ij}Y_{\ell m}(\omega),
\qquad 1\le i,j\le N.
\end{equation}

Second, for $|x|=r>R_j$, the standard partial-wave expansion of the free Green
function yields
\[
G_k^+(x,R_j\omega')
=
ik\sum_{n=0}^\infty \sum_{q=-n}^n
j_n(kR_j)\,h_n^{(1)}(kr)\,
Y_{nq}(\widehat x)\overline{Y_{nq}(\omega')}.
\]
Multiplying by $Y_{\ell m}(\omega')$ and integrating over $S^2$, orthonormality of
the spherical harmonics gives
\begin{equation}\label{eq:Phi-exterior-lemma}
\Phi_j(x,k)
=
ik\,\jj{\ell}(kR_j)\hj{\ell}(kr)\,Y_{\ell m}(\widehat x),
\qquad |x|=r>R_j.
\end{equation}
Similarly, for $|x|=r<R_j$, the same partial-wave expansion gives
\begin{equation}\label{eq:Phi-interior-lemma}
\Phi_j(x,k)
=
ik\,\hj{\ell}(kR_j)\,\jj{\ell}(kr)\,Y_{\ell m}(\widehat x),
\qquad |x|=r<R_j.
\end{equation}
Hence $\Phi_j(\cdot,k)$ lies in the fixed $(\ell,m)$ channel on each annulus.

Now set
\begin{equation}\label{eq:scatt-sol-lemma}
u_{\ell m}^+(x,k)
=
\jj{\ell}(k|x|)Y_{\ell m}(\widehat x)
-
\sum_{j=1}^N \theta_j c_{\ell,j}(k)\,\Phi_j(x,k),
\end{equation}
where
\[
\theta_j=\alpha_jR_j^2,
\qquad j=1,\ldots,N,
\]
so that
\[
\Theta=\diag(\theta_1,\ldots,\theta_N).
\]
By \eqref{eq:Phi-exterior-lemma} and \eqref{eq:Phi-interior-lemma},
the function $u_{\ell m}^+(\cdot,k)$ also lies in the fixed $(\ell,m)$
channel on each annulus.

Since $\jj{\ell}(k|x|)Y_{\ell m}(\widehat x)$ is regular at the origin and each
$\Phi_j(\cdot,k)$ is smooth near the origin, the function
$u_{\ell m}^+(\cdot,k)$ is regular at $x=0$. Moreover,
\[
(-\Delta-k^2)u_{\ell m}^+(\cdot,k)=0
\qquad \text{in }\R^3\setminus\bigcup_{j=1}^N S_j.
\]

We verify the interface conditions. Evaluating
\eqref{eq:scatt-sol-lemma} on $S_i$ and using
\eqref{eq:trace-Phi-lemma}, we get
\[
u_{\ell m}^+(R_i\omega,k)
=
\Bigl(
\jj{\ell}(kR_i)-\sum_{j=1}^N (m_\ell^+(k)\Theta)_{ij}c_{\ell,j}(k)
\Bigr)Y_{\ell m}(\omega).
\]
Since
\[
K_\ell^+(k)c_\ell(k)=b_\ell(k),
\]
this becomes
\[
u_{\ell m}^+(R_i\omega,k)=c_{\ell,i}(k)Y_{\ell m}(\omega).
\]

On the other hand, \eqref{eq:jump-Phi-lemma} gives
\begin{align*}
&\partial_r u_{\ell m}^+(R_i+0,\omega,k)
-
\partial_r u_{\ell m}^+(R_i-0,\omega,k)
\\
&\qquad
=
-\sum_{j=1}^N \theta_j c_{\ell,j}(k)
\bigl(
\partial_r\Phi_j(R_i+0,\omega,k)
-
\partial_r\Phi_j(R_i-0,\omega,k)
\bigr)
\\
&\qquad
=
\theta_i c_{\ell,i}(k)R_i^{-2}Y_{\ell m}(\omega)
=
\alpha_i c_{\ell,i}(k)Y_{\ell m}(\omega)
=
\alpha_i\,u_{\ell m}^+(R_i\omega,k).
\end{align*}
Thus $u_{\ell m}^+(\cdot,k)$ satisfies the transmission conditions for the
$\delta$-shell interaction and solves
\[
(H-k^2)u_{\ell m}^+(\cdot,k)=0
\]
in the sense of distribution.

We next determine its exterior form. For $|x|=r>R_N$,
\eqref{eq:Phi-exterior-lemma} yields
\[
u_{\ell m}^+(x,k)
=
\Bigl(
\jj{\ell}(kr)
-
ik\,\hj{\ell}(kr)\,{}^t b_\ell(k)\Theta c_\ell(k)
\Bigr)Y_{\ell m}(\widehat x).
\]
Using
\[
\jj{\ell}(kr)=\frac12\bigl(\hj{\ell}(kr)+h_\ell^{(2)}(kr)\bigr),
\]
we obtain
\[
u_{\ell m}^+(x,k)
=
\frac12
\Bigl(
h_\ell^{(2)}(kr)+\sigma_\ell(k)\hj{\ell}(kr)
\Bigr)
Y_{\ell m}(\widehat x),
\qquad r>R_N,
\]
where
\[
\sigma_\ell(k)
=
1-2ik\,{}^t b_\ell(k)\Theta c_\ell(k).
\]
Since
\[
c_\ell(k)=K_\ell^+(k)^{-1}b_\ell(k),
\]
this is exactly \eqref{eq:sigma-lemma}.
By construction, the exterior part involves only $h_\ell^{(1)}(kr)$,
hence the solution is outgoing in the sense of the Sommerfeld
radiation condition.
\end{proof}
%%%%%%%%%%%%
We show that the outgoing solution in the $(\ell,m)$ channel is unique.

\begin{lem}
\label{lem:no-homogeneous-outgoing}
Fix $\ell\ge0$, fix $m$ with $-\ell\le m\le\ell$, and let $k>0$.
Let $w$ be a function in the $(\ell,m)$ channel such that:

\begin{itemize}
\item[(i)]
$w$ is regular at the origin,

\item[(ii)]
\[
(-\Delta-k^2)w=0
\qquad \text{in }\R^3\setminus\bigcup_{j=1}^N S_j,
\]

\item[(iii)]
$w$ is continuous across each sphere $S_j$ and satisfies
\[
\partial_r w(R_j+0,\omega)-\partial_r w(R_j-0,\omega)
=
\alpha_j w(R_j,\omega),
\qquad j=1,\ldots,N,
\]

\item[(iv)]
$w$ is outgoing and has zero incident part, that is,
for $r>R_N$ one has
\[
w(x)=\gamma\,\hj{\ell}(kr)Y_{\ell m}(\widehat x)
\]
for some $\gamma\in\C$.
\end{itemize}

Then
\[
w\equiv0
\qquad \text{on }\R^3.
\]
\end{lem}

\begin{proof}
By assumption, for $r>R_N$,
$$
w(x)=\gamma\,\hj{\ell}(kr)Y_{\ell m}(\widehat x)
$$
for some $\gamma\in\C$.
We recall that, as $r\to\infty$,
$$
h_\ell^{(1)}(kr)
=
(-i)^{\ell+1}\frac{e^{ikr}}{kr}\bigl(1+O(r^{-1})\bigr),
$$
so that $h_\ell^{(1)}$ represents an outgoing spherical wave.

We show that $\gamma=0$.
Fix $r>R_N$ and decompose the ball $\{|x|<r\}$ into
the annuli determined by the shells.
Applying Green's identity on each annulus and summing over all annuli, we obtain
\begin{align*}
\int_{|x|=r}\overline{w}\,\partial_r w\,d\sigma
&=
\int_{|x|<r}\bigl(|\nabla w|^2-k^2|w|^2\bigr)\,dx
\\
&\quad
+
\sum_{j=1}^N
\int_{S_j}\overline{w}\,
\bigl(
\partial_r w(R_j+0)-\partial_r w(R_j-0)
\bigr)\,d\sigma_j.
\end{align*}
Here $d\sigma$ denotes the surface measure on the outer sphere $\{|x|=r\}$,
while $d\sigma_j$ denotes the surface measure on the shell $S_j$.
There is no contribution from $x=0$ because $w$ is regular at the origin.
Using
\[
\partial_r w(R_j+0)-\partial_r w(R_j-0)=\alpha_j w\upharpoonright_{S_j},
\qquad \alpha_j\in\R,
\]
we see that the right-hand side is real. Hence
\begin{equation}\label{eq:imagflux-zero-lemma}
\Im \int_{|x|=r}\overline{w}\,\partial_r w\,d\sigma =0.
\end{equation}

On the other hand, the asymptotics
\[
\hj{\ell}(kr)
=
(-i)^{\ell+1}\frac{e^{ikr}}{kr}\bigl(1+O(r^{-1})\bigr),
\]
and
\[
\partial_r\hj{\ell}(kr)
=
(-i)^{\ell+1}\frac{e^{ikr}}{kr}
\bigl(ik-r^{-1}+O(r^{-2})\bigr)
\]
as $r\to\infty$ imply
\[
\Im \int_{|x|=r}\overline{w}\,\partial_r w\,d\sigma
=
\frac{|\gamma|^2}{k}\,\|Y_{\ell m}\|_{L^2(S^2)}^2
+
O(r^{-1}).
\]
Letting $r\to\infty$ and using \eqref{eq:imagflux-zero-lemma}, we conclude
that $\gamma=0$. Therefore
\[
w(x)=0
\qquad \text{for }|x|>R_N.
\]

Since $w$ lies in the fixed $(\ell,m)$ channel, we may write
\[
w(x)=\frac{g(r)}{r}Y_{\ell m}(\widehat x)
\]
on each interval $(R_j,R_{j+1})$, where
\[
R_0:=0,
\qquad
R_{N+1}:=\infty.
\]
Then $g$ satisfies
\[
g''(r)+\left(k^2-\frac{\ell(\ell+1)}{r^2}\right)g(r)=0
\]
on each such interval.

From $w=0$ on $(R_N,\infty)$ we obtain
\[
g(R_N+0)=0,
\qquad
g'(R_N+0)=0.
\]
Since $w$ is continuous across $r=R_N$, we also have
\[
g(R_N-0)=g(R_N+0)=0.
\]
Moreover, the jump condition
\[
\partial_r w(R_N+0)-\partial_r w(R_N-0)=\alpha_N w(R_N)
\]
implies
\[
g'(R_N+0)-g'(R_N-0)=\alpha_N g(R_N)=0,
\]
and hence
\[
g'(R_N-0)=0.
\]
Therefore the Cauchy data of $g$ vanish at $R_N$ on $(R_{N-1},R_N)$, so
uniqueness for the ODE implies $g\equiv0$ there. Repeating the same argument
across the shells, we conclude that
\[
w\equiv0
\qquad \text{on }\R^3.
\]
\end{proof}

\begin{cor}
\label{cor:partial-wave-uniqueness}
Fix $\ell\ge0$, fix $m$ with $-\ell\le m\le\ell$, and let $k>0$.
Then there exists at most one outgoing solution in the $(\ell,m)$ channel
whose incident part is
\[
\jj{\ell}(k|x|)Y_{\ell m}(\widehat x).
\]
\end{cor}

\begin{proof}
Suppose that $u$ and $\widetilde u$ are two outgoing solutions in the
$(\ell,m)$ channel with the same incident part
\[
\jj{\ell}(k|x|)Y_{\ell m}(\widehat x).
\]
Then
\[
w:=u-\widetilde u
\]
is regular at the origin, satisfies
\[
(-\Delta-k^2)w=0
\qquad \text{in }\R^3\setminus\bigcup_{j=1}^N S_j,
\]
obeys the same interface conditions, and has zero incident part.
Hence $w$ satisfies all assumptions of
Lemma~\ref{lem:no-homogeneous-outgoing}. Therefore
\[
w\equiv0,
\]
and so
\[
u=\widetilde u.
\]
Thus the outgoing solution is unique.
\end{proof}
%%%
\begin{prop}
\label{prop:Kell-invertible-positive}
For each $\ell\ge0$ and each $k>0$, the matrix
\[
K_\ell(k^2+i0)
\]
is invertible. Equivalently,
\[
Z_\ell
:=
\{k>0:\det K_\ell(k^2+i0)=0\}
=
\varnothing.
\]
\end{prop}

\begin{proof}
Fix $\ell\ge0$, fix $k>0$, and choose $m$ with $-\ell\le m\le\ell$.
Write
\[
Y:=Y_{\ell m}.
\]
Assume, for contradiction, that $K_\ell(k^2+i0)$ is not invertible.
Then there exists a nonzero vector
\[
c={}^t(c_1,\ldots,c_N)\in\C^N
\]
such that
\[
K_\ell(k^2+i0)c=0.
\]

Let
\[
G_k^+(x,y):=\frac{e^{ik|x-y|}}{4\pi|x-y|},
\]
and for $j=1,\ldots,N$ define
\[
\Phi_j(x,k)
:=
\int_{S^2}G_k^+(x,R_j\omega')\,Y(\omega')\,d\omega'.
\]
As in the proof of Lemma~\ref{lem:partial-wave-construction}, each
$\Phi_j(\cdot,k)$ lies in the fixed $(\ell,m)$ channel, is regular at the
origin, satisfies
\[
(-\Delta-k^2)\Phi_j(\cdot,k)=0
\qquad
\text{in }\R^3\setminus S_j,
\]
is continuous across each sphere, and obeys
\[
\Phi_j(R_i\omega,k)
=
(m_\ell^+(k))_{ij}Y(\omega),
\qquad
1\le i,j\le N,
\]
together with
\[
\partial_r\Phi_j(R_i+0,\omega,k)-\partial_r\Phi_j(R_i-0,\omega,k)
=
-\delta_{ij}R_i^{-2}Y(\omega).
\]
Moreover, for $r>R_N$ one has
\[
\Phi_j(x,k)
=
ik\,j_\ell(kR_j)\,h_\ell^{(1)}(kr)\,Y(\widehat x),
\qquad |x|=r>R_N,
\]
since $R_j<R_N<r$.

Now define
\[
u(x)
=
-\sum_{j=1}^N (\Theta c)_j\,\Phi_j(x,k).
\]
Then $u$ lies in the fixed $(\ell,m)$ channel, is regular at the origin,
satisfies
\[
(-\Delta-k^2)u=0
\qquad
\text{in }\R^3\setminus\bigcup_{j=1}^N S_j,
\]
and is outgoing with zero incident part, because for $r>R_N$ it is a linear
combination of $h_\ell^{(1)}(kr)Y(\widehat x)$ only.

We check the interface conditions. For $1\le i\le N$,
\[
u(R_i\omega)
=
-\sum_{j=1}^N (m_\ell^+(k))_{ij}(\Theta c)_j\,Y(\omega)
=
-(m_\ell^+(k)\Theta c)_i\,Y(\omega).
\]
Since
\[
K_\ell(k^2+i0)c
=
(I_N+m_\ell^+(k)\Theta)c
=
0,
\]
we have
\[
m_\ell^+(k)\Theta c=-c,
\]
and therefore
\[
u(R_i\omega)=c_iY(\omega).
\]
Similarly,
\begin{align*}
&\partial_r u(R_i+0,\omega)-\partial_r u(R_i-0,\omega)
\\
&\qquad
=
-\sum_{j=1}^N (\Theta c)_j
\bigl(
\partial_r\Phi_j(R_i+0,\omega,k)-\partial_r\Phi_j(R_i-0,\omega,k)
\bigr)
\\
&\qquad
=
(\Theta c)_iR_i^{-2}Y(\omega)
=
\alpha_i c_iY(\omega)
=
\alpha_i u(R_i\omega).
\end{align*}
Thus $u$ satisfies the transmission conditions.

Since $c\neq0$, there exists some $i$ such that $c_i\neq0$, and hence
\[
u(R_i\omega)=c_iY(\omega)\not\equiv0.
\]
Therefore $u\not\equiv0$.

We have thus constructed a nontrivial outgoing solution in the
$(\ell,m)$ channel with zero incident part. This contradicts
Lemma~\ref{lem:no-homogeneous-outgoing}. Hence
$K_\ell(k^2+i0)$ must be invertible.

The final statement follows immediately.
\end{proof}

%%%
\subsection{Determination of the channel scattering coefficient}
We now identify the scattering coefficient in each channel and derive the determinant formula.

\begin{prop}
\label{prop:stationary-channel}
Fix $\ell \ge 0$ and $m$ with $-\ell \le m \le \ell$. Then, for almost every
$k>0$, the stationary scattering solution in the $(\ell,m)$ channel whose
incident part is
\[
j_\ell(k|x|)Y_{\ell m}(\widehat x)
=
\frac12\bigl(h_\ell^{(1)}(k|x|)+h_\ell^{(2)}(k|x|)\bigr)Y_{\ell m}(\widehat x)
\]
has exterior form
\begin{equation}\label{eq:stationary-channel}
u(x,k)=
\frac12\bigl(h_\ell^{(2)}(kr)+S_\ell(k)h_\ell^{(1)}(kr)\bigr)Y_{\ell m}(\widehat x),
\qquad r>R_N.
\end{equation}
Equivalently, if an outgoing solution in the $(\ell,m)$ channel with the same
incident part has exterior form
\[
\frac12\bigl(h_\ell^{(2)}(kr)+\beta h_\ell^{(1)}(kr)\bigr)Y_{\ell m}(\widehat x),
\qquad r>R_N,
\]
then $\beta=S_\ell(k)$.
\end{prop}

\begin{proof}
We use the standard stationary interpretation of the fiber scattering matrix
in the rotationally symmetric setting, with the present incoming/outgoing
normalization of the spherical waves.
For almost every $k>0$, the fiber operator $S(k^2)$ maps the incoming
amplitude at energy $k^2$ to the outgoing amplitude of the corresponding
stationary scattering solution, see
\cite[Ch.~XI, Sects.~8A--C]{ReedSimonIII} and \cite{Yafaev1992}.

Now
\[
j_\ell(k|x|)Y_{\ell m}(\widehat x)
=
\frac12 h_\ell^{(2)}(k|x|)Y_{\ell m}(\widehat x)
+
\frac12 h_\ell^{(1)}(k|x|)Y_{\ell m}(\widehat x),
\]
so, in the standard incoming/outgoing normalization for spherical waves, the
free channel wave has incoming amplitude $\frac12 Y_{\ell m}$.
Therefore the corresponding stationary scattering solution has exterior form
\[
\frac12 h_\ell^{(2)}(kr)Y_{\ell m}(\widehat x)
+
\frac12 h_\ell^{(1)}(kr)(S(k^2)Y_{\ell m})(\widehat x),
\qquad r>R_N.
\]

By Theorem~\ref{thm:wave},
\[
S(k^2)\upharpoonright_{{\cal H}_\ell}=S_\ell(k)I_{{\cal H}_\ell}
\]
for almost every $k>0$, hence
\[
S(k^2)Y_{\ell m}=S_\ell(k)Y_{\ell m}.
\]
Substituting this gives the stated formula.

The final assertion follows from this formula together with
Corollary~\ref{cor:partial-wave-uniqueness}.
\end{proof}
%%%


%%%%%%%%%%%%%%
We now state the main result of this section, which gives a determinant
formula for $S_\ell(k)$.

\begin{thm}
\label{thm:Sell}
For each $\ell\ge0$ and for almost every $k>0$,
\begin{equation}\label{eq:Sell-det}
S_\ell(k)
=
\frac{\det K_\ell(k^2-i0)}{\det K_\ell(k^2+i0)}.
\end{equation}
Moreover, by Proposition~\ref{prop:Kell-invertible-positive}, the
right-hand side of \eqref{eq:Sell-det} is well defined for every $k>0$.

In particular,
\[
|S_\ell(k)|=1
\qquad \text{for almost every }k>0,
\]
so one may choose a real-valued phase shift $\delta_\ell(k)$, defined for
almost every $k>0$, such that
\[
S_\ell(k)=e^{2i\delta_\ell(k)}.
\]
If $J\subset(0,\infty)$ is an interval, then, for any continuous branch of
$\arg \det K_\ell(k^2+i0)$ on $J$, one may choose $\delta_\ell$ on $J$ so that
\begin{equation}\label{eq:delta-arg}
\delta_\ell(k)
=
-\arg \det K_\ell(k^2+i0)
\qquad \text{for almost every }k\in J.
\end{equation}
\end{thm}

\begin{proof}
It suffices to prove \eqref{eq:Sell-det}.

Fix $\ell\ge0$, and let $k>0$ be such that $S_\ell(k)$ is defined.
By Theorem~\ref{thm:wave}, this holds for almost every $k>0$.
By Proposition~\ref{prop:Kell-invertible-positive},
\[
\det K_\ell(k^2+i0)\neq0.
\]

Fix $m$ with $-\ell\le m\le\ell$.
Hence Lemma~\ref{lem:partial-wave-construction} applies, and there exists
an outgoing solution in the $(\ell,m)$ channel with incident part
\[
\jj{\ell}(k|x|)Y_{\ell m}(\widehat x),
\]
and its exterior form is
\[
\frac12
\Bigl(
h_\ell^{(2)}(kr)+\sigma_\ell(k)\hj{\ell}(kr)
\Bigr)
Y_{\ell m}(\widehat x),
\qquad r>R_N,
\]
where
\[
\sigma_\ell(k)
=
1-2ik\,{}^t b_\ell(k)\Theta K_\ell(k^2+i0)^{-1}b_\ell(k).
\]
By Proposition~\ref{prop:stationary-channel}, the outgoing coefficient in
this normalization is exactly $S_\ell(k)$. Comparing
\eqref{eq:ext-form-lemma} with \eqref{eq:stationary-channel}, we obtain
\[
S_\ell(k)=\sigma_\ell(k).
\]
Therefore
\begin{equation}\label{eq:S-explicit-thm}
S_\ell(k)
=
1-2ik\,{}^t b_\ell(k)\Theta K_\ell(k^2+i0)^{-1}b_\ell(k).
\end{equation}

Moreover, using
\[
\hj{\ell}(t)+h_\ell^{(2)}(t)=2\jj{\ell}(t),
\]
we obtain
\[
(m_\ell(k^2-i0)-m_\ell(k^2+i0))_{ij}
=
-2ik\,\jj{\ell}(kR_i)\jj{\ell}(kR_j),
\qquad 1\le i,j\le N.
\]
Equivalently,
\[
m_\ell(k^2-i0)-m_\ell(k^2+i0)
=
-2ik\,b_\ell(k)\,{}^t b_\ell(k).
\]
Hence
\[
K_\ell(k^2-i0)
=
K_\ell(k^2+i0)-2ik\,b_\ell(k)\,{}^t b_\ell(k)\,\Theta.
\]

Applying the matrix determinant lemma
\[
\det(A+u\,{}^t v)=\det(A)\bigl(1+{}^t v A^{-1}u\bigr)
\]
with
\[
A=K_\ell(k^2+i0),
\qquad
u=-2ik\,b_\ell(k),
\qquad
{}^t v={}^t b_\ell(k)\Theta,
\]
we obtain
\begin{align*}
\det K_\ell(k^2-i0)
&=
\det K_\ell(k^2+i0)
\Bigl(
1-2ik\,{}^t b_\ell(k)\Theta K_\ell(k^2+i0)^{-1}b_\ell(k)
\Bigr)
\\
&=
\det K_\ell(k^2+i0)\,S_\ell(k)
\end{align*}
by \eqref{eq:S-explicit-thm}. Therefore
\[
S_\ell(k)
=
\frac{\det K_\ell(k^2-i0)}{\det K_\ell(k^2+i0)}.
\]
This proves \eqref{eq:Sell-det} for every $k>0$ such that $S_\ell(k)$ is
defined, hence for almost every $k>0$.

By Proposition~\ref{prop:Kell-invertible-positive}, the denominator in
\eqref{eq:Sell-det} is nonzero for every $k>0$, so the right-hand side is
well defined for every $k>0$.
By Lemma~\ref{lem:Kell-boundary},
\[
\det K_\ell(k^2-i0)
=
\overline{\det K_\ell(k^2+i0)}.
\]
Combining this with \eqref{eq:Sell-det}, we obtain
\[
|S_\ell(k)|=1
\qquad \text{for almost every }k>0.
\]
Hence one may choose a real-valued phase shift $\delta_\ell(k)$, defined for
almost every $k>0$, such that
\[
S_\ell(k)=e^{2i\delta_\ell(k)}.
\]

Now let $J\subset(0,\infty)$ be an interval, and set
\[
D_\ell(k):=\det K_\ell(k^2+i0).
\]
Since the entries of $K_\ell(k^2+i0)$ depend continuously on $k>0$,
the function $D_\ell$ is continuous on $(0,\infty)$.
By Proposition~\ref{prop:Kell-invertible-positive},
\[
D_\ell(k)\neq0
\qquad (k>0).
\]
Hence, on $J$ one may choose a continuous branch of $\arg D_\ell(k)$, and for
almost every $k\in J$ one has
\[
S_\ell(k)
=
\frac{\overline{D_\ell(k)}}{D_\ell(k)}
=
e^{-2i\arg D_\ell(k)}.
\]
Therefore one may choose $\delta_\ell$ on $J$ so that
\[
\delta_\ell(k)
=
-\arg D_\ell(k)
=
-\arg\det K_\ell(k^2+i0)
\qquad \text{for almost every }k\in J,
\]
which proves \eqref{eq:delta-arg}.
\end{proof}
%%%%%
\begin{rem}
By Proposition~\ref{prop:Kell-invertible-positive}, the determinant ratio in
\eqref{eq:Sell-det} is well defined for every $k>0$.
The phrase ``for almost every $k>0$'' refers only to the measurable
representative of the channel scattering coefficient arising from the
abstract scattering operator.
\end{rem}

%%%%%%
%%%%%%%%%%%%%%%%%%%%%%%%%%%%%%%%%%%%%%%%%%%%%%%%%%

%%%%%%%%%%%%%%%%%%%%%%%%%%%%%%%%%%%%%%%%%
\section{The double $\delta$--shell case}\label{sec:double}

We restrict attention to the $s$--wave channel ($\ell=0$), which captures
the leading contribution in the low--energy regime. In this channel the
formulas reduce to elementary functions and the threshold behavior can be
analyzed in a completely explicit form.

For $\ell\ge1$, the same determinant formula remains valid. In the present
paper, however, we restrict the threshold analysis to the $s$--wave channel,
which already captures the phenomenon of interest in the double--shell model.

We now specialize to the case $N=2$. Thus
$$
0<R_1<R_2,
\qquad
H=-\Delta+\alpha_1\delta(|x|-R_1)+\alpha_2\delta(|x|-R_2).
$$
We write
$$
\theta_j=\alpha_jR_j^2,
\qquad
j=1,2.
$$
In this case the general determinant formula from the previous section becomes
completely explicit. We restrict attention to the $s$--wave channel, where all
relevant quantities can be written in elementary functions.

\subsection{The $s$--wave channel}

We consider the case $\ell=0$. Then
$$
\jj{0}(x)=\frac{\sin x}{x},
\qquad
\hj{0}(x)=-\,\frac{i e^{ix}}{x}.
$$
For convenience, we also set
$$
s_j=\sin(kR_j),
\qquad
c_j=\cos(kR_j),
\qquad
j=1,2.
$$

The following lemma gives the boundary matrix in this channel.

\begin{lem}\label{lem:m0}
For $k>0$, the $s$--wave boundary matrix $m_0(k^2+i0)$ is given by
$$
m_0(k^2+i0)
=
\left(
\begin{array}{cc}
\dfrac{s_1e^{ikR_1}}{kR_1^2} &
\dfrac{s_1e^{ikR_2}}{kR_1R_2}
\\[12pt]
\dfrac{s_1e^{ikR_2}}{kR_1R_2} &
\dfrac{s_2e^{ikR_2}}{kR_2^2}
\end{array}
\right).
$$
Accordingly,
$$
K_0(k^2+i0)=I_2+m_0(k^2+i0)\Theta,
\qquad
\Theta=\diag(\theta_1,\theta_2).
$$
\end{lem}

\begin{proof}
This follows immediately from Lemma~\ref{lem:mell} with $\ell=0$, using
$$
j_0(t)=\frac{\sin t}{t},
\qquad
h_0^{(1)}(t)=-i\frac{e^{it}}{t}.
$$
The formula for $K_0(k^2+i0)$ is the definition of $K_\ell(z)$ specialized to
$\ell=0$.
\end{proof}

We now derive an explicit formula for $S_0(k)$ in the $s$--wave channel,
equivalently for the determinant $\det K_0(k^2+i0)$.

\begin{thm}
\label{thm:S0}
Define real functions $A_0(k)$ and $B_0(k)$ by
\begin{eqnarray}
A_0(k)
&=&
R_1^2R_2^2k^2
+
R_1^2k\,c_2s_2\,\theta_2
+
R_2^2k\,c_1s_1\,\theta_1
\nonumber\\
&&\hspace{1cm}
+
\theta_1\theta_2
\bigl(c_1c_2s_1s_2-c_2^2s_1^2\bigr),
\label{eq:A0}
\end{eqnarray}
and
\begin{eqnarray}
B_0(k)
&=&
R_1^2k\,s_2^2\,\theta_2
+
R_2^2k\,s_1^2\,\theta_1
\nonumber\\
&&\hspace{1cm}
+
\theta_1\theta_2
s_1s_2\bigl(c_1s_2-c_2s_1\bigr).
\label{eq:B0}
\end{eqnarray}
Then
\begin{equation}\label{eq:scaled-det}
k^2R_1^2R_2^2\,\det K_0(k^2+i0)=A_0(k)+iB_0(k).
\end{equation}
Consequently, for almost every $k>0$,
\begin{equation}\label{eq:S0}
S_0(k)=\frac{A_0(k)-iB_0(k)}{A_0(k)+iB_0(k)}.
\end{equation}
Moreover, for any interval $J\subset(0,\infty)$ and any continuous branch of
$\arg\bigl(A_0(k)+iB_0(k)\bigr)$ on $J$, one may choose the phase shift
$\delta_0$ on $J$ so that
\begin{equation}\label{eq:delta0}
\delta_0(k)
=
-\arg\bigl(A_0(k)+iB_0(k)\bigr),
\qquad \mbox{for almost every~} k\in J.
\end{equation}
\end{thm}

\begin{proof}
By Lemma~\ref{lem:m0},
$$
K_0(k^2+i0)
=
\left(
\begin{array}{cc}
1+\theta_1\dfrac{s_1e^{ikR_1}}{kR_1^2} &
\theta_2\dfrac{s_1e^{ikR_2}}{kR_1R_2}
\\[12pt]
\theta_1\dfrac{s_1e^{ikR_2}}{kR_1R_2} &
1+\theta_2\dfrac{s_2e^{ikR_2}}{kR_2^2}
\end{array}
\right).
$$
Hence
\begin{eqnarray}
\det K_0(k^2+i0)
&=&
\Bigl(1+\theta_1\dfrac{s_1e^{ikR_1}}{kR_1^2}\Bigr)
\Bigl(1+\theta_2\dfrac{s_2e^{ikR_2}}{kR_2^2}\Bigr)
\nonumber\\
&&\hspace{1cm}
-
\theta_1\theta_2
\dfrac{s_1^2e^{2ikR_2}}{k^2R_1^2R_2^2}.
\label{eq:detK0-step1}
\end{eqnarray}
Multiplying by $k^2R_1^2R_2^2$, we obtain
\begin{eqnarray}
k^2R_1^2R_2^2\,\det K_0(k^2+i0)
&=&
k^2R_1^2R_2^2
+
kR_2^2\theta_1 s_1e^{ikR_1}
+
kR_1^2\theta_2 s_2e^{ikR_2}
\nonumber\\
&&\hspace{1cm}
+
\theta_1\theta_2 s_1s_2e^{ik(R_1+R_2)}
-
\theta_1\theta_2 s_1^2e^{2ikR_2}.
\label{eq:detK0-step2}
\end{eqnarray}

Using
$$
e^{ikR_j}=c_j+is_j,
\qquad
e^{ik(R_1+R_2)}=(c_1+is_1)(c_2+is_2),
\qquad
e^{2ikR_2}=(c_2+is_2)^2,
$$
we separate the real and imaginary parts. 
A direct computation yields
\begin{equation}\label{eq:detK0-step3}
k^2R_1^2R_2^2\,\det K_0(k^2+i0)=A_0(k)+iB_0(k),
\end{equation}
where $A_0(k)$ and $B_0(k)$ are given by \eqref{eq:A0} and \eqref{eq:B0}.
This proves \eqref{eq:scaled-det}.

The formula for $S_0(k)$ now follows from Theorem~\ref{thm:Sell}. Indeed,
by \eqref{eq:scaled-det},
$$
\det K_0(k^2+i0)=\frac{A_0(k)+iB_0(k)}{k^2R_1^2R_2^2},
$$
and, by Lemma~\ref{lem:Kell-boundary},
\[
\det K_0(k^2-i0)
=
\overline{\det K_0(k^2+i0)}
=
\frac{A_0(k)-iB_0(k)}{k^2R_1^2R_2^2}.
\]
Substituting these expressions into \eqref{eq:Sell-det}, we obtain
\eqref{eq:S0} for almost every $k>0$.

Now let $J\subset(0,\infty)$ be an interval.
By Proposition~\ref{prop:Kell-invertible-positive},
\[
\det K_0(k^2+i0)\neq0
\qquad (k>0).
\]
Hence, by \eqref{eq:scaled-det},
\[
A_0(k)+iB_0(k)\neq0
\qquad (k>0).
\]
Therefore, for any continuous branch of
$\arg\bigl(A_0(k)+iB_0(k)\bigr)$ on $J$, one may choose $\delta_0$ on $J$ so
that
\[
\delta_0(k)
=
-\arg\det K_0(k^2+i0)
=
-\arg\bigl(A_0(k)+iB_0(k)\bigr),
\qquad \text{for almost every }k\in J,
\]
which is exactly \eqref{eq:delta0}.
\end{proof}

\begin{rem}
Formula \eqref{eq:S0} gives a completely explicit expression for the
$s$--wave scattering matrix in the double $\delta$--shell case.
\end{rem}
%%%%%%%%%%%%%%%%%%%%%%%%%%%%%%%%%%%%%%%%%%%%%%%%%%%
\section{Low--energy behavior in the double $\delta$--shell case}
\label{sec:low}

We continue to assume that $N=2$ and restrict attention to the $s$--wave
channel $\ell=0$. Using the explicit formula in
Theorem~\ref{thm:S0}, we analyze the behavior of the scattering matrix near
the threshold $k=0$. We first derive the regular threshold expansion and, in
particular, an explicit formula for the scattering length.

\begin{thm}
\label{thm:regular-threshold}
Set
\begin{equation}\label{eq:C0}
C_0
=
R_1^2R_2^2
+
R_1^2R_2\theta_2
+
R_1R_2^2\theta_1
+
\theta_1\theta_2R_1(R_2-R_1).
\end{equation}
Assume that
$$
C_0\neq0.
$$
Then one may choose the phase shift $\delta_0(k)$ for all sufficiently small
$k>0$ so that
$$
\delta_0(k)\to0
\qquad (k\downarrow0).
$$
For this choice, as $k\downarrow0$,
\begin{equation}\label{eq:delta0-low}
\delta_0(k)=-a_{\mathrm s}k+O(k^3),
\end{equation}
where the scattering length is given by
\begin{equation}\label{eq:scatt-length}
a_{\mathrm s}
=
\frac{
R_1^2R_2^2(\theta_1+\theta_2)
+
\theta_1\theta_2R_1R_2(R_2-R_1)
}{
R_1^2R_2^2
+
R_1^2R_2\theta_2
+
R_1R_2^2\theta_1
+
\theta_1\theta_2R_1(R_2-R_1)
}.
\end{equation}
Equivalently,
\begin{equation}\label{eq:S0-low}
S_0(k)=1-2ia_{\mathrm s}k+O(k^2),
\qquad
k\downarrow0.
\end{equation}
\end{thm}

\begin{proof}
We first show that
\[
A_0(k)+iB_0(k)\neq0
\]
for all sufficiently small $k>0$.
Once this is established, Theorem~\ref{thm:S0} allows us to choose the phase
shift so that
\[
\delta_0(k)=-\arg\bigl(A_0(k)+iB_0(k)\bigr)
\]
for all sufficiently small $k>0$.
We begin by expanding $A_0(k)$ and $B_0(k)$ as $k\downarrow0$.

Using
$$
\sin(kR_j)=kR_j+O(k^3),
\qquad
\cos(kR_j)=1+O(k^2),
\qquad
j=1,2,
$$
formula \eqref{eq:A0} gives
\begin{eqnarray}
A_0(k)
&=&
R_1^2R_2^2k^2
+
R_1^2k\,\theta_2\bigl(kR_2+O(k^3)\bigr)
+
R_2^2k\,\theta_1\bigl(kR_1+O(k^3)\bigr)
\nonumber\\
&&\hspace{1cm}
+
\theta_1\theta_2\bigl(R_1R_2-R_1^2\bigr)k^2
+
O(k^4).
\label{eq:A0-low-expand1}
\end{eqnarray}
Hence
\begin{equation}\label{eq:A0-low}
A_0(k)=C_0\,k^2+O(k^4).
\end{equation}

Similarly, \eqref{eq:B0} yields
\begin{eqnarray}
B_0(k)
&=&
R_1^2k\,\theta_2\bigl(k^2R_2^2+O(k^4)\bigr)
+
R_2^2k\,\theta_1\bigl(k^2R_1^2+O(k^4)\bigr)
\nonumber\\
&&\hspace{1cm}
+
\theta_1\theta_2
(kR_1)(kR_2)\bigl(kR_2-kR_1\bigr)
+
O(k^5),
\label{eq:B0-low-expand1}
\end{eqnarray}
and therefore
\begin{equation}\label{eq:B0-low}
B_0(k)=\Gamma_0\,k^3+O(k^5),
\end{equation}
where
\begin{equation}\label{eq:Gamma0}
\Gamma_0
=
R_1^2R_2^2(\theta_1+\theta_2)
+
\theta_1\theta_2R_1R_2(R_2-R_1).
\end{equation}

Since $C_0\neq0$, relation \eqref{eq:A0-low} shows that
$$
A_0(k)\neq0
$$
for all sufficiently small $k>0$. 
Hence also
$$
A_0(k)+iB_0(k)\neq0
$$
for all sufficiently small $k>0$. Moreover, by \eqref{eq:S0},
$$
S_0(k)
=
\frac{1-i\,B_0(k)/A_0(k)}{1+i\,B_0(k)/A_0(k)}.
$$
Since
$$
\frac{B_0(k)}{A_0(k)}=O(k)
\qquad (k\downarrow0),
$$
it follows that
$$
S_0(k)=1+O(k)
\qquad (k\downarrow0).
$$
Therefore, one may choose the phase shift $\delta_0(k)$ for all sufficiently
small $k>0$ so that
$$
\delta_0(k)\to0
\qquad (k\downarrow0).
$$
Since phase shifts are determined only modulo $\pi$, for this choice we have
\begin{equation}\label{eq:delta-arctan}
\delta_0(k)
=
-\arctan\!\frac{B_0(k)}{A_0(k)}
\end{equation}
for all sufficiently small $k>0$.

Moreover,
$$
\frac{B_0(k)}{A_0(k)}=O(k)
\qquad (k\downarrow0).
$$
Using
$$
\arctan t=t+O(t^3)
\qquad (t\to0),
$$
we obtain from \eqref{eq:delta-arctan} that
\begin{equation}\label{eq:arg-expand}
\delta_0(k)
=
-\frac{B_0(k)}{A_0(k)}+O(k^3).
\end{equation}
Combining \eqref{eq:A0-low} and \eqref{eq:B0-low}, we obtain
\begin{equation}\label{eq:delta0-low-proof}
\delta_0(k)
=
-\frac{\Gamma_0}{C_0}\,k+O(k^3).
\end{equation}
Thus \eqref{eq:delta0-low} holds with
$$
a_{\mathrm s}=\frac{\Gamma_0}{C_0},
$$
which is exactly \eqref{eq:scatt-length}.

Finally, since
$$
S_0(k)=e^{2i\delta_0(k)},
$$
equation \eqref{eq:delta0-low} implies
$$
S_0(k)=1+2i\delta_0(k)+O(k^2)
=
1-2ia_{\mathrm s}k+O(k^2),
$$
which proves \eqref{eq:S0-low}.
\end{proof}


%%%%%%%%%%%%%%%%%%%% new
\begin{rem}
The condition $C_0\neq0$ characterizes the regular threshold regime.
The complementary case $C_0=0$ corresponds to a threshold--critical
configuration.

Under the standing assumption $0<R_1<R_2$, the condition $C_0=0$ already
implies
\[
\Gamma_0\neq0,
\]
as will be shown in the proof of Theorem~\ref{thm:exceptional-threshold}.
More precisely, the next theorem treats the nondegenerate exceptional case in
which
\[
C_0=0,
\qquad
C_2\neq0,
\]
where $C_2$ is defined in Theorem~\ref{thm:exceptional-threshold}.
Further degenerate situations, such as $C_0=0$ and $C_2=0$, require a
separate higher--order analysis.
\end{rem}

We next consider the nondegenerate exceptional threshold regime in the
double $\delta$--shell case.

\begin{thm}
\label{thm:exceptional-threshold}
Assume that $N=2$ and $\ell=0$, and define $C_0$ and $\Gamma_0$ by
\eqref{eq:C0} and \eqref{eq:Gamma0}. Suppose that
\[
C_0=0,
\qquad
C_2\neq0,
\]
where
\begin{align}
C_2
&=
-\frac23 R_1^2R_2^3\theta_2
-\frac23 R_1^3R_2^2\theta_1
\nonumber\\
&\quad
+
\theta_1\theta_2
\left[
-\frac23(R_1^3R_2+R_1R_2^3)
+R_1^2R_2^2+\frac13 R_1^4
\right].
\label{eq:C2}
\end{align}
Then, as $k\downarrow0$,
\begin{equation}\label{eq:A0-exceptional}
A_0(k)=C_2 k^4+O(k^6),
\end{equation}
and
\begin{equation}\label{eq:B0-exceptional}
B_0(k)=\Gamma_0 k^3+O(k^5).
\end{equation}
In particular,
\begin{equation}\label{eq:BA-exceptional}
\frac{B_0(k)}{A_0(k)}
=
\frac{\Gamma_0}{C_2}\frac{1}{k}
+O(k),
\qquad k\downarrow0,
\end{equation}
and hence
\begin{equation}\label{eq:S0-exceptional}
S_0(k)\to -1
\qquad (k\downarrow0).
\end{equation}
Equivalently, one may choose the phase shift $\delta_0(k)$ on
$(0,\varepsilon)$, for some $\varepsilon>0$, so that
\begin{equation}\label{eq:delta-exceptional}
\delta_0(k)\to \pm \frac{\pi}{2}
\qquad (k\downarrow0).
\end{equation}
\end{thm}

\begin{proof}
We first note that, under the standing assumption $0<R_1<R_2$, the condition
$C_0=0$ already implies
\[
\Gamma_0\neq0.
\]
Indeed, if $\Gamma_0=0$ as well, then
\[
0
=
\frac{C_0}{R_1}-\frac{\Gamma_0}{R_1R_2}
=
R_2\bigl(R_1R_2+(R_2-R_1)\theta_1\bigr),
\]
and hence
\[
R_1R_2+(R_2-R_1)\theta_1=0.
\]
Substituting this into
\begin{eqnarray*}
\Gamma_0
&=&
R_1R_2\Bigl(
R_1R_2(\theta_1+\theta_2)+(R_2-R_1)\theta_1\theta_2
\Bigr) \\
&=&
R_1R_2\Bigl(
\theta_2\bigl(R_1R_2+(R_2-R_1)\theta_1\bigr)+R_1R_2\theta_1
\Bigr),
\end{eqnarray*}
we obtain
\[
\Gamma_0=R_1^2R_2^2\theta_1
=
-\frac{R_1^3R_2^3}{R_2-R_1}\neq0,
\]
a contradiction.

We next expand $A_0(k)$ to higher order as $k\downarrow0$.

Using
\[
\sin(kR_j)=kR_j-\frac{(kR_j)^3}{6}+O(k^5),
\qquad
\cos(kR_j)=1-\frac{(kR_j)^2}{2}+O(k^4),
\qquad j=1,2,
\]
we first compute
\[
c_js_j
=
\left(1-\frac{k^2R_j^2}{2}+O(k^4)\right)
\left(kR_j-\frac{k^3R_j^3}{6}+O(k^5)\right)
=
kR_j-\frac{2}{3}k^3R_j^3+O(k^5).
\]
Hence
\begin{align*}
R_1^2k\,c_2s_2\,\theta_2
&=
R_1^2R_2\theta_2\,k^2
-\frac{2}{3}R_1^2R_2^3\theta_2\,k^4
+O(k^6),
\\
R_2^2k\,c_1s_1\,\theta_1
&=
R_1R_2^2\theta_1\,k^2
-\frac{2}{3}R_1^3R_2^2\theta_1\,k^4
+O(k^6).
\end{align*}

Next, we expand the last term in \eqref{eq:A0}. We have
\begin{align*}
s_1s_2
&=
k^2R_1R_2
-\frac{k^4}{6}(R_1^3R_2+R_1R_2^3)
+O(k^6),
\\
c_1c_2
&=
1-\frac{k^2}{2}(R_1^2+R_2^2)+O(k^4),
\end{align*}
and therefore
\begin{align*}
c_1c_2s_1s_2
&=
k^2R_1R_2
-\frac{2}{3}k^4(R_1^3R_2+R_1R_2^3)
+O(k^6).
\end{align*}
Also,
\begin{align*}
s_1^2
&=
k^2R_1^2-\frac{1}{3}k^4R_1^4+O(k^6),
\\
c_2^2
&=
1-k^2R_2^2+O(k^4),
\end{align*}
so
\begin{align*}
c_2^2s_1^2
&=
k^2R_1^2-k^4\left(R_1^2R_2^2+\frac13 R_1^4\right)+O(k^6).
\end{align*}
Thus
\begin{align*}
c_1c_2s_1s_2-c_2^2s_1^2
&=
k^2R_1(R_2-R_1)
\\
&\quad
+
k^4\left[
-\frac23(R_1^3R_2+R_1R_2^3)
+R_1^2R_2^2+\frac13 R_1^4
\right]
+O(k^6).
\end{align*}

Substituting these expansions into \eqref{eq:A0}, we obtain
\[
A_0(k)=C_0k^2+C_2k^4+O(k^6).
\]
Since $C_0=0$, this gives \eqref{eq:A0-exceptional}.

On the other hand, by \eqref{eq:B0-low},
\[
B_0(k)=\Gamma_0k^3+O(k^5),
\]
which proves \eqref{eq:B0-exceptional}.

Since $C_2\neq0$, we may divide \eqref{eq:B0-exceptional} by
\eqref{eq:A0-exceptional} and obtain
\[
\frac{B_0(k)}{A_0(k)}
=
\frac{\Gamma_0}{C_2}\frac{1}{k}+O(k),
\qquad k\downarrow0.
\]
Because $C_0=0$ implies $\Gamma_0\neq0$, the coefficient $\Gamma_0/C_2$ is
nonzero. Hence
\[
\left|\frac{B_0(k)}{A_0(k)}\right|\to\infty
\qquad (k\downarrow0).
\]
Using \eqref{eq:S0}, we write
\[
S_0(k)
=
\frac{1-i\,B_0(k)/A_0(k)}{1+i\,B_0(k)/A_0(k)}.
\]
Since $B_0(k)/A_0(k)\to\pm\infty$, it follows that
\[
S_0(k)\to -1
\qquad (k\downarrow0),
\]
which proves \eqref{eq:S0-exceptional}.

Since $C_0=0$ implies $\Gamma_0\neq0$, relation \eqref{eq:B0-exceptional}
shows that
\[
B_0(k)\neq0
\]
for all sufficiently small $k>0$. Hence
\[
A_0(k)+iB_0(k)\neq0
\qquad (0<k<\varepsilon)
\]
for some $\varepsilon>0$. Since $A_0(k)+iB_0(k)$ is continuous and
nonvanishing on $(0,\varepsilon)$, one may choose a continuous branch of
$\arg\bigl(A_0(k)+iB_0(k)\bigr)$ there. Therefore, by
Theorem~\ref{thm:S0}, one may choose the phase shift on $(0,\varepsilon)$ so
that
\[
\delta_0(k)=-\arg\bigl(A_0(k)+iB_0(k)\bigr).
\]
In view of \eqref{eq:S0-exceptional}, this choice satisfies
\[
\delta_0(k)\to \pm\frac{\pi}{2}
\qquad (k\downarrow0),
\]
which proves \eqref{eq:delta-exceptional}.
\end{proof}
We show that the condition $C_0=0$ is equivalent to the existence of a
nontrivial zero--energy radial solution whose exterior constant term
vanishes.

\begin{prop}
\label{prop:C0-zero-energy}
Assume that $N=2$ and $\ell=0$. Then the following are equivalent:

\begin{itemize}
\item[(i)]
\[
C_0=0.
\]

\item[(ii)]
There exists a nontrivial radial function $u$, piecewise harmonic away from
the shells, continuous across the shells, and satisfying the $\delta$--shell
jump conditions, such that
\[
Hu=0
\]
in the distributional sense, $u$ is regular at the origin, and
\[
u(x)=O(|x|^{-1})
\qquad (|x|\to\infty).
\]
\end{itemize}

More precisely, every radial zero--energy solution regular at the origin is of
the form
\[
u(x)=f(|x|),
\]
where
\[
f(r)=
\begin{cases}
a, & 0<r<R_1,\\[2mm]
b+\dfrac{c}{r}, & R_1<r<R_2,\\[3mm]
d+\dfrac{e}{r}, & r>R_2,
\end{cases}
\]
with constants satisfying
\[
c=-\theta_1 a,
\qquad
b=a+\frac{\theta_1}{R_1}a,
\]
and
\[
d=\frac{C_0}{R_1^2R_2^2}\,a.
\]
Hence the exterior constant term vanishes if and only if $C_0=0$.
\end{prop}

\begin{proof}
Let $u(x)=f(r)$ be a radial solution of $Hu=0$, where $r=|x|$. Away from the
shells, the equation reduces to
\[
-\Delta u=0.
\]
For an $s$--wave radial function in three dimensions, the general harmonic
solution is of the form
\[
f(r)=A+\frac{B}{r}.
\]
Regularity at the origin forces $B=0$ in the region $0<r<R_1$. Hence
\[
f(r)=
\begin{cases}
a, & 0<r<R_1,\\[2mm]
b+\dfrac{c}{r}, & R_1<r<R_2,\\[3mm]
d+\dfrac{e}{r}, & r>R_2,
\end{cases}
\]
for suitable constants $a,b,c,d,e$.

We impose continuity and the $\delta$--shell jump conditions.

At $r=R_1$, continuity gives
\[
a=b+\frac{c}{R_1}.
\]
Since
\[
f'(r)=0 \quad (0<r<R_1),
\qquad
f'(r)=-\frac{c}{r^2}\quad (R_1<r<R_2),
\]
the jump condition at $r=R_1$ yields
\[
-\frac{c}{R_1^2}
=
\alpha_1 a.
\]
Using $\theta_1=\alpha_1R_1^2$, we obtain
\[
c=-\theta_1 a,
\qquad
b=a+\frac{\theta_1}{R_1}a.
\]

At $r=R_2$, continuity gives
\[
b+\frac{c}{R_2}=d+\frac{e}{R_2}.
\]
Moreover,
\[
f'(r)=-\frac{c}{r^2}\quad (R_1<r<R_2),
\qquad
f'(r)=-\frac{e}{r^2}\quad (r>R_2),
\]
so the jump condition at $r=R_2$ becomes
\[
-\frac{e}{R_2^2}+\frac{c}{R_2^2}
=
\alpha_2\left(b+\frac{c}{R_2}\right).
\]
Equivalently,
\[
e
=
c-\theta_2\left(b+\frac{c}{R_2}\right).
\]
Substituting the expressions for $b$ and $c$, we obtain
\[
e
=
-\theta_1 a
-
\theta_2
\left(
a+\frac{\theta_1}{R_1}a-\frac{\theta_1}{R_2}a
\right).
\]

It remains to compute the exterior constant term $d$. From continuity at
$r=R_2$,
\[
d
=
b+\frac{c}{R_2}-\frac{e}{R_2}.
\]
Substituting the formulas above and simplifying, we find
\[
d
=
\left(
1+\frac{\theta_1}{R_1}
+\frac{\theta_2}{R_2}
+\theta_1\theta_2
\left(
\frac{1}{R_1R_2}-\frac{1}{R_2^2}
\right)
\right)a.
\]
Multiplying by $R_1^2R_2^2$, this becomes
\[
R_1^2R_2^2\,d
=
\Bigl(
R_1^2R_2^2
+
R_1R_2^2\theta_1
+
R_1^2R_2\theta_2
+
\theta_1\theta_2R_1(R_2-R_1)
\Bigr)a
=
C_0a.
\]
Hence
\[
d=\frac{C_0}{R_1^2R_2^2}\,a.
\]

Therefore the exterior constant term vanishes if and only if $d=0$, that is,
if and only if $C_0=0$.

Conversely, any piecewise harmonic radial function that is continuous across
$r=R_1,R_2$ and satisfies the above jump conditions is a distributional
solution of $Hu=0$. Moreover, if $a=0$, then the formulas obtained above give
\[
b=c=d=e=0,
\]
so the solution is trivial. Hence a nontrivial radial distributional solution,
regular at the origin and satisfying
\[
u(x)=O(|x|^{-1})
\qquad (|x|\to\infty),
\]
exists if and only if $a\neq0$ and $d=0$, which is equivalent to $C_0=0$.
\end{proof}

We show that the scattering length is determined by the exterior behavior of
the zero--energy radial solution.

\begin{prop}
\label{prop:scatt-length-zero-energy}
Assume that $N=2$, $\ell=0$, and $C_0\neq0$. Let $u$ be a nontrivial radial
solution of
\[
Hu=0
\]
which is regular at the origin, and normalize it so that
\[
u(x)=1-\frac{a}{|x|}
\qquad (|x|>R_2).
\]
Then
\[
a=a_{\mathrm s},
\]
where $a_{\mathrm s}$ is given by \eqref{eq:scatt-length}.
\end{prop}

\begin{proof}
In the notation of Proposition~\ref{prop:C0-zero-energy}, the exterior part of
a radial zero--energy solution is
\[
d+\frac{e}{r},
\]
and the computation in the proof of
Proposition~\ref{prop:C0-zero-energy} gives
\[
d=\frac{C_0}{R_1^2R_2^2}a_0,
\qquad
e=-\frac{\Gamma_0}{R_1^2R_2^2}a_0,
\]
where $a_0$ denotes the interior constant on $(0,R_1)$. Since $C_0\neq0$, we
may normalize by $d=1$. Then
\[
u(x)
=
1+\frac{e}{|x|}
=
1-\frac{\Gamma_0}{C_0}\frac{1}{|x|},
\qquad |x|>R_2.
\]
By \eqref{eq:scatt-length}, one has
\[
a_{\mathrm s}=\frac{\Gamma_0}{C_0}.
\]
Therefore
\[
a=a_{\mathrm s}.
\]
\end{proof}
%%%%%%%%%%%%%%%%%%%%%%%%%
Proposition~\ref{prop:C0-zero-energy} gives a concrete interpretation of the
exceptional threshold condition in
Theorem~\ref{thm:exceptional-threshold}. In the regular case $C_0\neq0$, a
zero--energy radial solution regular at the origin has a nonzero constant term
at infinity, and the standard scattering--length description applies. By
contrast, the condition $C_0=0$ means that the exterior constant term
vanishes, leaving a decaying tail of order $r^{-1}$.

Thus the exceptional regime corresponds to a threshold--critical configuration
in which the contributions of the two shells cancel at zero energy. In this
case,
$$
S_0(k)\to -1
\qquad (k\downarrow0)
$$
as shown in Theorem~\ref{thm:exceptional-threshold}. 
This behavior reflects the presence of a zero--energy $s$--wave solution with vanishing exterior
constant term and explains the breakdown of the finite scattering--length
picture.

\begin{rem}
A further degenerate situation may occur if $C_0=0$ and $C_2=0$. In that case,
higher--order terms in the threshold expansion become relevant and require a
separate analysis.
\end{rem}

\section*{Concluding Remarks}

In this paper we derived a determinant formula for the channel scattering
coefficients of Schr\"odinger operators with finitely many concentric
$\delta$--shell interactions. The main result shows that, after partial--wave
reduction, the scattering problem is reduced to a finite-dimensional matrix
problem governed by the same boundary operator that appears in the resolvent
formula.

As an application, we analyzed in detail the double--shell model in the
$s$--wave channel. We obtained explicit formulas for the scattering matrix and
the scattering length, and we gave a zero--energy characterization of a
threshold--critical configuration. In particular, we showed that the condition
$C_0=0$ is equivalent to the existence of a zero--energy radial solution with
vanishing exterior constant term. In the corresponding nondegenerate
exceptional case, this threshold--critical configuration yields the limiting
behavior
$$
S_0(k)\to -1
\qquad (k\downarrow0).
$$

The determinant formula also admits a natural structural interpretation.
Whenever
$$
\|m_\ell(z)\Theta\|<1,
$$
where $\|\cdot\|$ denotes the operator norm of matrices on $\C^N$,
the inverse of the reduced boundary matrix is given by the convergent
Neumann series
$$
K_\ell(z)^{-1}
=
(I_N+m_\ell(z)\Theta)^{-1}
=
I_N-m_\ell(z)\Theta+(m_\ell(z)\Theta)^2-\cdots.
$$
Each term in this series represents one additional step in the multiple
scattering process between the shells.
Thus $K_\ell(z)$ provides a finite-dimensional description of this process,
and the determinant in the scattering formula may be viewed as encoding the
cumulative effect of these repeated interactions. This interpretation is
therefore not merely formal in regimes where the above Neumann series
converges, for example for sufficiently small coupling strengths.

Several natural problems remain for further study. One important problem is
to relate the determinant formula to the spectral shift function and the
Birman--Kre\u{\i}n formula \cite{BirmanKrein1962}. Another is
to analyze more degenerate threshold situations, where higher-order terms in
the expansion become relevant. It would also be interesting to extend the
present approach to more general hypersurface interactions.

\small
\begin{thebibliography}{99}

\bibitem{AlbeverioEtAl2005}
S.~Albeverio, F.~Gesztesy, R.~H{\o}egh--Krohn, and H.~Holden,
Solvable Models in Quantum Mechanics,
2nd ed.,
AMS Chelsea Publishing,
Providence, RI, 2005.
DOI: 10.1090/chel/350

\bibitem{BehrndtLangerLotoreichik2013}
J.~Behrndt, M.~Langer, and V.~Lotoreichik,
Schr\"odinger operators with $\delta$ and $\delta'$-potentials supported on hypersurfaces,
Ann.\ Henri Poincar\'e 14 (2013), 385--423.
DOI: 10.1007/s00023-012-0189-5

\bibitem{BirmanKrein1962}
M.~Sh.~Birman and M.~G.~Kre\u{\i}n,
On the theory of wave operators and scattering operators,
Dokl.\ Akad.\ Nauk SSSR 144 (1962), 475--478;
English transl.: Soviet Math.\ Dokl.\ 3 (1962), 740--744.

\bibitem{BrascheEtAl1994}
J.~F.~Brasche, P.~Exner, Yu.~A.~Kuperin, and P.~{\v S}eba,
Schr\"odinger operators with singular interactions,
J.\ Math.\ Anal.\ Appl.\ 184 (1994), 112--139.
DOI: 10.1006/jmaa.1994.1188

\bibitem{HounkonnouHounkpeShabani1997}
M.~N.~Hounkonnou, M.~Hounkpe, and J.~Shabani,
Scattering theory for finitely many sphere interactions supported by concentric spheres,
J.\ Math.\ Phys.\ 38 (1997), 2832--2850.
DOI: 10.1063/1.532022

\bibitem{IkebeShimada1991}
T.~Ikebe and S.~Shimada,
Spectral and scattering theory for the Schr\"odinger operators with
penetrable wall potentials,
J.\ Math.\ Kyoto Univ.\ 31 (1991), 219--258.
DOI: 10.1215/kjm/1250519902

\bibitem{KaminagaAAMP2026}
M.~Kaminaga,
Schr\"odinger operators with concentric $\delta$--shell interactions,
to appear in Analysis and Mathematical Physics; arXiv:2602.09376.

\bibitem{Posilicano2001}
A.~Posilicano,
A Krein-like formula for singular perturbations of self-adjoint operators and applications,
J.\ Funct.\ Anal.\ 183 (2001), 109--147.
DOI: 10.1006/jfan.2000.3730

\bibitem{ReedSimonIII}
M.~Reed and B.~Simon,
Methods of Modern Mathematical Physics. III: Scattering Theory,
Academic Press,
New York, 1979.

\bibitem{Shabani1988}
J.~Shabani,
Finitely many $\delta$ interactions with supports on concentric spheres,
J.\ Math.\ Phys.\ 29 (1988), 660--664.
DOI: 10.1063/1.528005

\bibitem{Yafaev1992}
D.~R.~Yafaev,
Mathematical Scattering Theory: General Theory,
Translations of Mathematical Monographs, Vol.~105,
American Mathematical Society,
Providence, RI, 1992.
DOI: 10.1090/mmono/105

\end{thebibliography}
\end{document}
